% Zone „Das kennst du schon“ – aus bank/winkel-dreiecke/zone.jsonl (zusammenbau v0.1)
\einheitenkopf[zone][Kennst du schon]{Das kennst du schon}
\setcounter{aufgabe}{0}
%% TODO Grundvorstellungs-Aufgabe als erste Hauptnummer der Zone (2.8) fehlt in der Bank

%% TODO Titel: Ich-kann-Satz fehlt in der Bank (Platzhalter: Kettenname) – Nr. 1
%% TODO Anweisung: ein Satz über der ersten Teilaufgabe fehlt in der Bank – Nr. 1
\begin{aufgabe}{Mit Lineal und Geodreieck zeichnen}
\swz{Miss mit dem Lineal die Länge und die Breite des Rechtecks.}{\rechteck{5}{3}}
\swa{Zeichne eine Gerade g und einen Punkt P, der 3 cm von g entfernt liegt. Zeichne mit dem Geodreieck die Senkrechte zu g durch P und die Parallele zu g durch P.}{\begin{ksys}[xmin=0,xmax=10,ymin=0,ymax=6] \end{ksys}}
\end{aufgabe}

%% TODO Titel: Ich-kann-Satz fehlt in der Bank (Platzhalter: Kettenname) – Nr. 2
%% TODO Anweisung: ein Satz über der ersten Teilaufgabe fehlt in der Bank – Nr. 2
\begin{aufgabe}{Addieren und Subtrahieren im Bereich bis zum Vollwinkel, auch mit Dezimalzahlen (Vielfaches minus zwei Dezimalgrade)}
%% TODO Darstellung für schwach fehlt in der Bank (winkel-dreiecke-zone-f2-v1; Abschnitt „Für schwache Schüler“)
\swz{$38^\circ + 29^\circ$}{}
%% TODO Darstellung für schwach fehlt in der Bank (winkel-dreiecke-zone-f2-v3; Abschnitt „Für schwache Schüler“)
\swz{$96{,}4^\circ + 38{,}7^\circ$}{}
\end{aufgabe}

%% TODO Titel: Ich-kann-Satz fehlt in der Bank (Platzhalter: Kettenname) – Nr. 3
%% TODO Anweisung: ein Satz über der ersten Teilaufgabe fehlt in der Bank – Nr. 3
\begin{aufgabe}{Längeneinheiten m und km umrechnen (1,5 km = 1500 m)}
%% TODO Darstellung für schwach fehlt in der Bank (winkel-dreiecke-zone-f3-v1; Abschnitt „Für schwache Schüler“)
\swz{3 km in Meter}{}
%% TODO Darstellung für schwach fehlt in der Bank (winkel-dreiecke-zone-f3-v3; Abschnitt „Für schwache Schüler“)
\swz{2,4 km in Meter}{}
\end{aufgabe}

%% TODO Titel: Ich-kann-Satz fehlt in der Bank (Platzhalter: Kettenname) – Nr. 4
%% TODO Anweisung: ein Satz über der ersten Teilaufgabe fehlt in der Bank – Nr. 4
\begin{aufgabe}{Vierecksarten kennen}
%% TODO Darstellung für schwach fehlt in der Bank (winkel-dreiecke-zone-f4-v1; Abschnitt „Für schwache Schüler“)
\swz[3]{Wie heißt ein Viereck mit vier rechten Winkeln und vier gleich langen Seiten?}{}
\swz[3]{Welche Seiten des Vierecks sind parallel? Wie heißt das Viereck?}{\trapez{5}{3}{2.5}}
\end{aufgabe}

%% TODO Titel: Ich-kann-Satz fehlt in der Bank (Platzhalter: Kettenname) – Nr. 5
%% TODO Anweisung: ein Satz über der ersten Teilaufgabe fehlt in der Bank – Nr. 5
\begin{aufgabe}{Kreis mit dem Zirkel zeichnen, Radius und Durchmesser unterscheiden}
\swa{Zeichne mit dem Zirkel einen Kreis mit dem Radius 3 cm.}{\begin{ksys}[xmin=0,xmax=10,ymin=0,ymax=6] \end{ksys}}
\swz{Welche Strecke ist der Radius r, welche der Durchmesser d? Miss beide.}{\begin{kreis}[2] \mittelpunkt{M} \radius{40}{r} \durchmesser{0}{d} \end{kreis}}
\end{aufgabe}

%% TODO Titel: Ich-kann-Satz fehlt in der Bank (Platzhalter: Kettenname) – Nr. 6
%% TODO Anweisung: ein Satz über der ersten Teilaufgabe fehlt in der Bank – Nr. 6
\begin{aufgabe}{Addieren und Subtrahieren im Bereich bis zum Vollwinkel, auch mit Dezimalzahlen (Vielfaches minus zwei Dezimalgrade) – Fehler finden}
\swz[3]{Mia rechnet: \rechnung{180^\circ - 46{,}8^\circ - 37{,}5^\circ &= 96{,}3^\circ} Finde den Fehler und rechne richtig.}{}
\end{aufgabe}

%% TODO Titel: Ich-kann-Satz fehlt in der Bank (Platzhalter: Kettenname) – Nr. 7
%% TODO Anweisung: ein Satz über der ersten Teilaufgabe fehlt in der Bank – Nr. 7
\begin{aufgabe}{Addieren und Subtrahieren im Bereich bis zum Vollwinkel, auch mit Dezimalzahlen (Vielfaches minus zwei Dezimalgrade) – selbst rechnen}
%% TODO Darstellung für schwach fehlt in der Bank (winkel-dreiecke-zone-f2-v6; Abschnitt „Für schwache Schüler“)
\swz{$180^\circ - 38{,}6^\circ - 55{,}7^\circ$}{}
\end{aufgabe}
